\section{A path-containment assertion}
\label{app:opposition}

Theorem~53 (pp.~16--17) of version~1 of Gonz\'alez Nevado's
\cite{gonzalez-nevado-2026} states a positive solution of the geometric
Generalized Lax conjecture. Its proof invokes Theorem~52 (p.~16), which asserts
containment of a fixed rigidly convex set throughout a smooth real-zero
deformation from a product of normalized tangent linear factors to a
multiple of the defining polynomial. The following example violates
that intermediate assertion under its stated hypotheses.

A real polynomial $g$ with $g(0)=1$ is \emph{real-zero} if the roots of
$s\mapsto g(sx)$ are real for every real vector $x$. Its closed
rigidly convex set is the closure of the component of $\{g\ne0\}$
containing the origin. Fix $a,b,\eta>0$ and, in one variable, put
\[
 g(x)=(1-x/a)(1+x/b),\qquad
 c(t)=1+\eta\sin^2(\pi t),\qquad
 H_t(x)=g(c(t)x),\quad 0\le t\le1.
\]
The zero set of $g$ is the compact smooth set $\{-b,a\}$.
The normalized tangent linear factors at these two points are
$1+x/b$ and $1-x/a$, whose product is $g$. Every $H_t$ has value $1$
at the origin and two distinct real roots,
\[
 -\frac{b}{c(t)}\quad\hbox{and}\quad\frac{a}{c(t)}.
\]
Thus the path depends smoothly on $t$ and consists of real-zero
polynomials with compact smooth zero sets. It even admits the smooth
monic real symmetric determinantal representation
\[
 H_t(x)=\det\begin{pmatrix}
 1-c(t)x/a&0\\
 0&1+c(t)x/b
 \end{pmatrix}.
\]
Both endpoints equal $g$, so the final polynomial is $g$ times the
constant cofactor $1$, as allowed by the stated hypotheses.
Nevertheless the rigidly convex set of $H_t$ is
\[
 \left[-\frac{b}{c(t)},\frac{a}{c(t)}\right],
\]
which is strictly smaller than $[-b,a]$ whenever $0<t<1$.
Consequently it does not contain the rigidly convex set of $g$.

The parameters $a,b>0$ describe the open family of normalized real
quadratics with negative leading coefficient, so the failure is not
confined to a symmetric or multiple-root example. The parameter $\eta$
may be arbitrarily small. The cited Theorem~52 states no restriction
to two or more variables, no lower bound on pencil size, and no
requirement that the endpoint cofactor be nonconstant.

This example refutes the all-intermediate-times conclusion of that
theorem. Its endpoints coincide, so it does not refute an endpoint-only
statement. The nonspectrahedrality result of this paper follows from
the construction and obstruction proved above.
